\chapter{Customer Partnerships}
\label{cap:customerpartnerships}


Customer Partnership represents an important and strategic pillar to accomplish the Bank's goals. The IT Department must establish measurable relationships between external customers and other stakeholders, assuming a role that actively contributes to the improvement of the services provided and, in specific, monitors through KPIs the value perceived by customers.
The IT has to collaborate also internally within the organization in order to engage the appropriate department whenever it is needed.



This topic is fundamental for the success of the Bank's vision, because the international expansion through acquisitions and other partnerships introduces new business requirements due to different customer expectations, cultural diversity, and increased service complexity.
These challenges must be supported in a proactive way, by ensuring that IT services remain aligned with business objectives, deliver consistent value across different countries, and contribute to the Bank's growth, while complying with agreed Service Level Agreements between customers, stakeholders and user expectations.




These requisites, makes performance measurement crucial. By continuously monitoring and collecting customer-oriented metrics and analyzing services scores, IT Department is able to implement a continual service improvements that can strengthen customer trust and enrich the service quality. Gaining long-term partneship.
This customer-centric approach supports organization's objectives of delivering high-performance, high-quality and cost-effective IT services. Aimed toward to accomplish the predefined goals \ref{cap:goals}


\section{Time To Market}
%   Savas

Time to market, also known as Speed to Market is made operational through \textbf{Average Lead Time for Change to Production} KPI metric, defines as the elapsed time between formal approval of a business requirement and its deployment into the production environment. This metric is selected because it directly reflects the organization’s ability to translate customer and business demands into operational banking capabilities which is essential for achieving the stated rapid digital service evolution and revenue growth objectives of the organisation. In the context of the strategic goal of the organisation is to increase the income from internet banking, reduced Time to Market provides higher organizational agility and faster delivery of customer targeting features, thus improving market responsiveness and competitive positioning.

There are several ITIL aligned ITSM processes that primarily effects Time to Market. \textbf{Change enablement} has a direct effect through approval workflows and Risk Management procedures, were excessive governance overhead can result in significant delays. \textbf{Release and Deployment Management} also effects the indicator  as deployment scheduling , release batching and automation maturity determine how quickly the validated changes end up deployed in the production. \textbf{Service Validation and Testing} influences Lead Time through test cycle duration and error handling loops, which can create bottlenecks if automation coverage is insufficient. \textbf{Service Design and Transition} activities also effect indirectly throught the time required for requirement analysis, architectural validation, and coordination of technical dependencies. Overall, inefficiencies or lack of cooperation across these processes will result in the increase of Lead Time, where higher process integration, standardization and automation would reduce it, resulting with the improvement of Time To Market KPI.



\section{Performance To Contract}
PerformanceToContract --> A